\scanpage{238}
- \mspage{186} - \margin{IX}
\[
{}_n A^0\simeq\underline R\otimes_{\mathbb Z}\mathbb Z/n\mathbb Z\ ,\tag{13.5.5}
\]
qui donne, pour $n$ parcourant les puissances d'un nombre premier $\ell$ donné, un isomorphisme de systèmes projectifs de groupes sur $S$
\[
T_\ell(A^0)\simeq\underline R\otimes_{\mathbb Z}\mathbb Z_\ell\ ,\tag{13.5.6}
\]
où $\mathbb Z_\ell$ est considéré comme le système projectif des $\mathbb Z/\ell^{n+1}\mathbb Z$. Dans le cas envisagé ici où $A$ provient d'un schéma en groupes séparé sur $S$, on voit alors immédiatement que les $\underline R\otimes\mathbb Z/n\mathbb Z$ peuvent être considérés comme des schémas relatifs, ce qui donne un sens au symbole $(\underline R\otimes_{\mathbb Z}\mathbb Z_\ell)_{\underline{\unsure{O|0}}}$ (alors que $R_{\underline{\unsure{O|0}}}$ n'en aurait pas en général !), et on obtient un isomorphisme canonique
\[
T_\ell(A_{\underline{\unsure{O|0}}})\unsure{\overset{?}{\simeq}|\simeq}(\underline R\otimes_{\mathbb Z}\mathbb Z_\ell)_{\underline{\unsure{O|0}}}\ .\tag{13.5.7}
\]
Il résulte alors immédiatement des constructions du par. 2 et de 13.3 que l'isomorphisme précédent transforme partie fixe en partie fixe, partie torique en partie torique :
\[
T_\ell(A_{\underline{\unsure{O|0}}})^{\mathrm f}\simeq(\underline R^{\mathrm f}\otimes_{\mathbb Z}\mathbb Z_\ell)_{\underline{\unsure{O|0}}}\ ,\quad T_\ell(A_{\underline{\unsure{O|0}}})^{\mathrm t}\simeq(\underline R^{\mathrm t}\otimes_{\mathbb Z}\mathbb Z_\ell)_{\underline{\unsure{O|0}}}\ .\tag{13.5.8}
\]
Si $A'$ est un groupe analytique relatif sur $S$ satisfaisant aux mêmes hypothèses que $A$, et si $W_{\underline{\unsure{O|0}}}$ est une biextension de $(A_{\underline{\unsure{O|0}}},A'_{\underline{\unsure{O|0}}})$ par $\mathbb G_{m\underline{\unsure{O|0}}}$, définissant donc au-dessus d'un voisinage de $s$ une biextension de $(A,A')$ par $\mathbb G_{mS}$ (en travaillant maintenant avec la notion au sens analytique complexe, obtenue en travaillant dans le site des espaces analytiques au-dessus de $S$, pour une topologie convenable \ldots), alors d'après VIII 2.2 $W$ définit des accouplements
\margin{504}
\margin{237}
\margin{504}
